\origpage{355}{N. Robertson, P.D. Seymour / Journal of Combinatorial Theory, Series B 90 (2004) 325--385}
\noindent obtained by cutting $\Sigma_\chi$ along $U(F) \cap \Sigma_\chi$; let $\Sigma'$ be the fragment with $\psi(\Sigma') \subseteq \Sigma_2$, where $\psi$ is the associated surjection. Let $\chi'$ be the colour scheme with $\Sigma_{\chi'} = \Sigma'$, $\Phi_{\chi'} = \psi^{-1}(\Phi_\chi) \cap \Sigma'$, $\Omega_{\chi'}(k) = \Omega_\chi(k)$ ($k = 2, 3$), and $\Omega_{\chi'}(S') = \Omega_\chi(S)$ for each long side $S'$ of $\Phi_{\chi'}$, where $S$ is the long side of $\Phi_\chi$ with $\psi(S') \subseteq S$. Since $\hat{\Sigma}_\chi$ is not a sphere, and $\hat{\Sigma}_{\chi'}$ is a sphere, and $\chi$ satisfies $\mathbf{S}_1$, it follows that $\chi'$ is not bad. Now for each $i \ge 1$, let $\Gamma_i' = \psi^{-1}(\Gamma_i) \cap \Sigma'$, and for each $e \in E(\Gamma_i')$, let $\phi_i'(e) = \phi_i(\psi(e))$ if $e$ is not a short side of $\Phi_{\chi'}$, and $\phi_i'(e) = e$ if $e$ is a short side. Then $(\Gamma_i', \phi_i')$ is a $\chi'$-coloured painting. Since $\chi'$ is not bad, there exist $j > i \ge 1$ and a linear inflation $\sigma'$ of $(\Gamma_i', \phi_i')$ in $(\Gamma_j', \phi_j')$. For each $e \in E(\Gamma_j \cap \Sigma_2)$, let $\sigma(e) = \psi(\sigma'(\psi^{-1}(e)))$, and for each $v \in V(\Gamma_j \cap \Sigma_2)$, let $\sigma(v) = \psi(\sigma'(v'))$, where $\psi^{-1}(v) = \{v'\}$. Then $\sigma$ is an inflation satisfying the three conditions above, as required. This proves 10.2. $\square$

\textbf{10.3.} \textit{Let $\chi$ satisfy $\mathbf{S}_1$, and let $n \ge 1$ be an integer. Then there is a well-behaved set $\mathscr{C}_1(n)$ of $\chi$-places with the following property. Let $(\Gamma, \phi, \lambda, \mathscr{T}) \in \mathscr{D}$ satisfy $rep(\Gamma), ord(\mathscr{T}) \ge n > rep(\mathscr{T})$, and $dist(\Gamma) \ge \frac{1}{2}n + 1$. Then there is a rooted location $\mathscr{L}$ which $n$-isolates $\mathscr{T}$ and for which $(\Gamma, \phi, \mathscr{L}) \in \mathscr{C}_1(n)$.}

\textbf{Proof.} By 5.2, there are finitely many pairs $(F_i, \pi_i)$ ($i \in I$), such that each $F_i$ is an even length circuit in $\hat{\Sigma}_\chi$ with $\frac{1}{2}|E(F_i)| < n$, and $\pi_i$ is a march with $\bar{\pi}_i \subseteq V(F_i)$ consisting of every second vertex of $F_i$, with the following property. Let $(\Gamma, \phi, \lambda, \mathscr{T})$ be as in the theorem, and let $F$ be a representativeness flaw. Then there exist $i \in I$ and a $\Phi_\chi$-preserving homeomorphism $\alpha$ of $\Sigma_\chi$ which maps $F$ to $F_i$ and $U(F) \cap V(\Gamma)$ to $\bar{\pi}_i$. Let $\mathscr{C}_1(n)$ be the union, over all $i \in I$, of the set of all images of members of $\mathscr{C}(F_i, \pi_i)$ (defined in 10.2) under $\Phi_\chi$-preserving homeomorphisms. By 9.4, $\mathscr{C}_1(n)$ is well-behaved (since $I$ is finite). Let $(\Gamma, \phi, \lambda, \mathscr{T}), F, i, \alpha$ be as before, and let $(\Gamma', \phi', \lambda', \mathscr{T}')$ be the image of $(\Gamma, \phi, \lambda, \mathscr{T})$ under $\alpha$ (in the natural sense). Since $\mathscr{S}$ is closed under $\Phi_\chi$-preserving homeomorphisms, it follows that $(\Gamma', \phi', \lambda', \mathscr{T}') \in \mathscr{D}$, and $F_i$ is a representativeness flaw for it with $U(F_i) \cap V(\Gamma') = \bar{\pi}_i$. By 10.2, there is a rooted location $\mathscr{L}'$ with $ord(\mathscr{L}') = \frac{1}{2}|E(F)| < n$ which $n$-isolates $\mathscr{T}'$ and with $(\Gamma', \phi', \mathscr{L}') \in \mathscr{C}(F_i, \pi_i)$. Let $\mathscr{L}$ be the image of $\mathscr{L}'$ under $\alpha^{-1}$. Then it satisfies the theorem. $\square$

The sets $\mathscr{C}_1(n)$ will handle failure of representativeness. Now we turn to another possible failure when $rep(\mathscr{T})$ is large but $d(r_1^*, r_2^*)$ is small for two poles $r_1^*, r_2^*$. More precisely, for $n \ge 1$ an integer, let us say that $(\Gamma, \phi, \lambda, \mathscr{T}) \in \mathscr{D}$ is \textit{$n$-flawed in distance} if $rep(\mathscr{T}) > n$, and $dist(\Gamma) > \frac{1}{2}n + 1$, and there are distinct poles $r_1^*, r_2^*$ with $d(r_1^*, r_2^*) \le n$. Suppose that $(\Gamma, \phi, \lambda, \mathscr{T}) \in \mathscr{D}$ is $n$-flawed in distance. Then since $dist(\Gamma) \ge \frac{1}{2}n + 1$, it follows from Theorem 9.2 of [6] that there exists $F \subseteq \Gamma^*$ such that one of the following holds:

\hspace*{12pt}$\bullet$ $F$ is a circuit with $|E(F)| \le 2n$, satisfying

\hspace*{24pt}$\circ$ either $|U(F) \cap bd(\Sigma_\chi)| \le 1$ and $U(F) \subseteq \Sigma_\chi$, or $|U(F) \cap bd(\Sigma_\chi)| = 2$ and $U(F) \not\subseteq \Sigma_\chi$, and
